\addtotoc{Хураангуй}

\vspace{10mm}

\begin{center}
\textbf{\large Хураангуй}
\end{center}

\sloppy

Энэхүү судалгаа нь санхүүгийн цаг хугацааны цуваан дээр (EUR/USD) машин сургалтын ансамбль ашиглан гурван ангиллын (BUY / HOLD / SELL) дохиог автоматаар гаргах системийг боловсруулсан. Гадаад валютын зах зээлийн (Форекс) шуугиан ихтэй, тогтворгүй (non-stationary) шинж чанар нь ангиллын хүндрэлтэй сорилт болдог. Загварын ангиллын гүйцэтгэлийг цаг хугацааны дарааллыг хадгалсан баталгаажуулалт болон бодит нөхцөлийг дуурайсан симуляцаар цогцоор үнэлэв.

LightGBM, XGBoost, CatBoost гурван GBDT загварыг зөөлөн санал хураалтаар (soft voting) нэгтгэсэн ансамблийг EUR/USD-ийн 2015--2024 оны өгөгдөл дээр сургасан. Урагш гүйлгэн баталгаажуулалтын (Walk-Forward Validation, WFV) Fold~5 огтлолын 2024 оны тестийн ерөнхий нарийвчлал 89.25\% (F1-macro=59.45\%, F1-weighted=87.23\%) бөгөөд HOLD ангийн давамгайлсан тэнцвэргүй өгөгдөлд нарийвчлалыг ганц хэмжүүр болгох хангалтгүй гэдгийг F1-macro харуулна. Баталгаажуулалтын өгөгдөл дээр итгэлцэл $\geq$0.90 дохионы нарийвчлал 96.96\% гарсан нь босгон шүүлтүүрийн үндэслэлийг харуулна. MetaTrader~5 дээр 2025 оны бүтэн жилд бүх 10 итгэлцлийн сценари эерэг өгөөж үзүүлсэн; 0.91 босгон дээр +517.70\% (\$61{,}770) хүрч, суурь жишгийн (EUR/USD-г эхнээс эцэс хүртэл барих стратеги) +13.47\%-г мэдэгдэхүйц давсан. Энэ дүн нь 1\% эрсдэлийн тогтмол хэмжээлэлт ба нийлмэл хуримтлагдах өсөлтийн үр нөлөөнөөс үүдэлтэй учир босго бүрийн сценарийг харьцуулах хэмжүүр болгон тайлбарлах нь зүйтэй. Брокерын шимтгэл болон шөнийн хадгалалтын зардлыг тусгаагүй зэрэг хязгаарлалтыг §\ref{sec:limitations}-аас үзнэ үү.

Системийн инженерчлэлийн талд офлайн загвар сургалт, REST API сервер (Flask + MongoDB), мобайл харагдах хэсэг (React Native), push мэдэгдэл (Expo Push) гэсэн дөрвөн давхаргыг нэгтгэсэн ML serving дамжлагыг ажиллах прототип хэлбэрээр бэлдсэн.

\vspace{3mm}
\textit{Түлхүүр үгс:} машин сургалт, GBDT ансамбль загвар, урагш гүйлгэн баталгаажуулалт, олон хугацааны интервалын шинж чанар, түүхэн өгөгдөл дээрх системийн туршилт, дохио-арилжааны хөрвүүлэлт, хүлээгдэж буй өгөөж ($R$), эрсдэл-өгөөжийн үнэлгээ
